\chapter{Hysteroscopy}

\section*{Introduction \& Clinical Indications}
Hysteroscopy uses a camera passed through the cervix to inspect and treat the uterine cavity. It can diagnose or remove endometrial polyps, submucosal fibroids, adhesions, retained tissue, selected septa, and focal lesions causing bleeding or infertility. It may be performed in an office or operating room according to the lesion, patient tolerance, comorbidity, equipment, and required intervention.

\section*{Relevant Surgical Anatomy}
The hysteroscope traverses the vagina, cervical canal, internal os, and uterine cavity. The endometrium, tubal ostia, anterior and posterior walls, fundus, and cervix are inspected. Fluid distension permits visualization, but excessive absorption can cause electrolyte disturbance or fluid overload. The myometrium and uterine serosa are at risk during deep resection.

\section*{Preoperative Workup \& Preparation}
Pregnancy is reasonably excluded. Assessment includes bleeding pattern, fertility plans, medications, anemia, infection, cervical stenosis, uterine imaging, and anesthesia needs. The lesion's size, depth, and location determine whether office or operative hysteroscopy is suitable. Consent covers pain, bleeding, infection, perforation, fluid overload, cervical injury, incomplete treatment, adhesions, and repeat procedures.

\section*{Surgical Technique \& Procedure}
The hysteroscope is introduced through the cervix, often without a speculum in office vaginoscopy. Saline or another approved distension medium expands the cavity. Biopsy forceps, scissors, electrosurgery, or a mechanical tissue-removal device may treat the lesion. Fluid input and output are monitored throughout. Specimens are sent for histopathology.

\section*{Complications \& Risks}
Risks include cramping, vasovagal symptoms, bleeding, infection, cervical laceration, uterine perforation, thermal injury, fluid overload, electrolyte abnormality, incomplete removal, and intrauterine adhesions. A suspected perforation or excessive fluid deficit requires immediate procedural reassessment.

\section*{Postoperative Care \& Recovery}
Most patients have light bleeding and cramping for a short period. Instructions cover analgesia, activity, intercourse, pathology, and warning signs. Fever, severe pain, heavy bleeding, fainting, foul discharge, or persistent vomiting requires review. Follow-up depends on pathology and the treated lesion.

\section*{Outcomes \& Prognosis}
Hysteroscopy provides direct visualization and targeted treatment, but symptoms can recur when the underlying hormonal or structural disease persists. Fertility and bleeding outcomes depend on lesion type, completeness of treatment, uterine factors, and patient goals. No universal success rate applies to all hysteroscopic procedures.

\section*{References}
\begin{enumerate}
  \item American College of Obstetricians and Gynecologists. The Use of Hysteroscopy for Diagnosis and Treatment of Intrauterine Pathology.
  \item American Association of Gynecologic Laparoscopists. Practice Guidelines for Hysteroscopy.
  \item Royal College of Obstetricians and Gynaecologists. Outpatient Hysteroscopy Guidance.
\end{enumerate}

\clearpage
